\newtheorem{defi}{Definition }[section]
\newtheorem{exemple}{Exemple}[section]
\newtheorem{rema}{Remarque}

\newcommand{\ttt}[1]{\texttt{#1}}

\newcommand{\nbCordes}{\ttt{nc}}
\newcommand{\nbFrettes}{\ttt{nf}}
\newcommand{\scordature}{\ttt{scorda}}

\newcommand{\tailleMain}{\ttt{tailleMain}}
\newcommand{\nbDoigtsUtil}{\ttt{nbDoigtsUtilisés}}
\newcommand{\isBarre}{\ttt{barréPossible}}
\newcommand{\nbPouce}{\ttt{nbCordesBloquéesPouce}}
\newcommand{\dLegale}{\ttt{distanceLégale}}

\newcommand{\accord}{\ttt{accord}}
\newcommand{\nbAccord}{\ttt{nbAccord}}
\newcommand{\lenAccord}{\ttt{lenAccord}}
\newcommand{\indetPremier}{\mathcal{F}}

\newcommand{\posiPhyPos}{\ttt{pos}\phi}

\newcommand{\indet}{\mathcal{X}}
\newcommand{\domaine}{\mathcal{D}}
\newcommand{\contr}{\mathcal{C}}
\newcommand{\modele}{\mathcal{M}}
\newcommand{\defmodele}{\modele = \left(\indet, \domaine, \contr\right)}
\newcommand{\ssi}{$\textit{ssi }$}


\newcommand\restr[2]{{% we make the whole thing an ordinary symbol
  \left.\kern-\nulldelimiterspace % automatically resize the bar with \right
  #1 % the function
  \littletaller % pretend it's a little taller at normal size
  \right|_{#2} % this is the delimiter
  }}

\newcommand{\littletaller}{\mathchoice{\vphantom{\big|}}{}{}{}}

\newcommand{\accolades}[1]{
  \left\{#1\right\}
}
\newcommand{\internat}[1]{
  \llbracket #1 \rrbracket
}
% écriture d'un accord sous la forme : fondamentale, liste de shifts
\newcommand{\chord}[2]{
    (#1\foreach \p in {#2}{, \p})
}


\newcommand{\tablature}[4]{  % représentation d'une tablature
%on a en arguments : échelle, nc, nf, positions
  \begin{tikzpicture}[scale=#1]
    % dessin du manche
    \draw[black, thick] (-0.5,1.5*#3) -- (#2-0.5,1.5*#3);
    \draw[black, thick] (-0.5,1.5*0.1+1.5*#3) -- (#2-0.5,1.5*0.1+1.5*#3);
    \newcounter{x}
    \forloop{x}{0}{\value{x} < #2}{
      \draw[gray, thick] (\value{x},0) -- (\value{x},1.5*#3);
    }

    \newcounter{y}
    \forloop{y}{0}{\value{y} < #3}{
      \draw[gray, thick] (-0.1,1.5*\value{y}) -- (#2-0.9,1.5*\value{y});
    }
    
    \newcounter{i}
    \setcounter{i}{0}
    \foreach \p in {#4}{
      \ifthenelse{\equal{\p}{0}}
        {\filldraw[green] (\value{i}, 1.5*#3) circle (9pt);}
        {\ifthenelse{\equal{\p}{-1}}
          {\filldraw[red] (\value{i}, 1.5*#3) circle (9pt);}
          {\filldraw[blue] (\value{i}, 0.75+1.5*#3-1.5*\p) circle (6pt);}
        }
      
      \addtocounter{i}{1}
    }
	\end{tikzpicture}
}


\newcommand{\code}[2]{  % pour écrire du code
  \begin{lstlisting}[language=#1]
      #2
	\end{lstlisting}
}



